\chapter{Business Intelligence, Looker, and Real-Time Serving}
\index{Looker}\index{LookML}\index{semantic layer}\index{reverse ETL}\index{Looker Studio}\index{explores}

\begin{objectives}
\begin{itemize}
    \item Explain Looker's architecture and the role of the LookML semantic layer in governed BI.
    \item Model dimensions, measures, views, and explores that encode metric definitions once.
    \item Contrast Looker Studio with Looker (Core) for self-serve versus governed analytics.
    \item Implement reverse ETL: push curated BigQuery segments to operational APIs with identity resolution.
    \item Build a real-time streaming-sales dashboard with drill-down filters and a reverse-ETL sync.
    \item Design a governed BI access layer for five hundred analysts with non-conflicting metrics.
    \item Complete the Milestone~4 capstone: an end-to-end real-time fraud-detection pipeline.
\end{itemize}
\end{objectives}

\begin{crosscloud}
Every platform eventually confronts the same BI governance question: are metrics defined once in a semantic layer, or independently in every dashboard?

\vspace{2mm}
{\footnotesize
\begin{tabular}{@{}p{2.8cm}p{3.0cm}p{3.0cm}p{3.0cm}@{}} \\
\toprule
\textbf{Capability} & \textbf{GCP Looker} & \textbf{AWS} & \textbf{Azure / Fabric} \\
\midrule
Governed BI + semantic layer & Looker (LookML) & QuickSight (datasets) & Power BI (semantic models) \\
Self-serve dashboards & Looker Studio & QuickSight & Power BI / Fabric \\
Semantic model as code & LookML in Git & QuickSight APIs & TMDL / Fabric git integration \\
In-memory acceleration & BI Engine (Ch. 6) & QuickSight SPICE & Import mode / DirectLake \\
Reverse ETL & Scheduled queries + Functions (partner tools) & Partner tools (Census, Hightouch) & Partner tools \\
\bottomrule
\end{tabular}}

\vspace{1mm}
\textbf{Snowflake} users pair with the same BI tools over its warehouses; \textbf{MongoDB} charts serve operational analytics directly on Atlas. The portable insight: a semantic layer is a \emph{code artifact}---versioned, reviewed, tested---and the organizations that treat it as such are the ones where ``Gross Revenue'' means one thing.
\end{crosscloud}

\begin{keyterms}
\textbf{LookML}: Looker's modeling language defining dimensions, measures, joins, and explores as versioned code.
\textbf{View}: a LookML file mapping a table or derived table into dimensions and measures.
\textbf{Explore}: a curated join of views exposed to analysts as a self-serve query surface.
\textbf{Measure}: a LookML aggregate definition (count, sum, average) with declared type and filters.
\textbf{Reverse ETL}: syncing curated warehouse data back into operational systems (CRM, marketing platforms).
\textbf{Identity resolution}: reconciling entity keys across systems so a segment in BigQuery matches the same humans in the CRM.
\end{keyterms}

\section{Week 16 Roadmap}
Week~16 closes Part~IV where streaming data meets human decision-making. Monday covers Looker's architecture and LookML. Tuesday models dimensions, measures, and explores. Wednesday contrasts self-serve and governed BI and introduces reverse ETL. Thursday builds the streaming-sales dashboard and a reverse-ETL sync to a marketing webhook. Friday designs the governed access layer for five hundred analysts. The week ends with the Milestone~4 real-time fraud capstone.

\begin{definitionbox}
\textbf{Looker} is Google Cloud's governed BI platform. Its defining feature is the \textbf{LookML semantic layer}: metric and join definitions live as versioned code between the database and every dashboard, so all consumers compute the same numbers \cite{googlecloudlookerdocs}.
\end{definitionbox}

\section{Monday: Looker Architecture and the Semantic Layer}
\index{Looker!architecture}\index{LookML!semantic layer}

\subsection{How Looker Differs}
Most BI tools extract data into extracts or connect dashboards directly to tables; Looker connects to the database (BigQuery, via a service account) and generates SQL per user interaction. Nothing is extracted by default; freshness is the warehouse's freshness. The price of this elegance is that warehouse performance (Chapters 5--6) \emph{is} dashboard performance---partitioning, clustering, and BI Engine decisions surface directly in tile load times.

\subsection{Why a Semantic Layer}
Without a shared layer, every analyst re-implements ``Gross Revenue'' in a SQL snippet, a sheet formula, or a dashboard calculated field, and the definitions diverge silently. LookML defines it once---\texttt{measure: gross\_revenue \{ type: sum sql: \$\{order\_amount\} ;; filters: [is\_refunded: "no"] \}}---and every explore, dashboard, and API consumer inherits the same definition, reviewable in Git like any other code.

\begin{notebox}
The semantic layer is also a security layer: LookML access grants and user attributes restrict which explores, fields, and rows a user can see, composing with (not replacing) the warehouse's row access policies from Chapter~8.
\end{notebox}

\section{Tuesday: Dimensions, Measures, Explores, and Views}
\index{dimensions}\index{measures}\index{explores!design}

\subsection{Modeling Vocabulary}
A \textbf{view} maps one table (or derived table) to LookML: \textbf{dimension}s for attributes and \textbf{measure}s for aggregates, each with types, SQL expressions, and formatting. An \textbf{explore} joins views with declared relationships and join keys, producing the analyst-facing query surface. Discipline points: always declare \texttt{relationship} (one\_to\_many, many\_to\_one) so Looker can apply symmetric aggregates and avoid fan-out inflation; expose only fields analysts should see; name measures as the business names them.

\begin{lstlisting}[language=SQL,caption={A minimal LookML view and explore (LookML syntax in SQL listing).}]
# file: views/order_items.view.lkml
view: order_items {
  sql_table_name: analytics.fact_order_items ;;
  dimension: order_id    { type: string sql: ${TABLE}.order_id ;; }
  dimension: order_date  { type: date   sql: ${TABLE}.order_date ;; }
  dimension: is_refunded { type: yesno  sql: ${TABLE}.refund_flag = 'Y' ;; }
  measure: gross_revenue {
    type: sum
    sql: ${TABLE}.order_amount ;;
    filters: [is_refunded: "no"]
    value_format_name: usd
  }
}

# file: models/retail.model.lkml
explore: orders {
  join: customers {
    sql_on: ${orders.customer_id} = ${customers.customer_id} ;;
    relationship: many_to_one
  }
}
\end{lstlisting}

\begin{pitfall}
Fan-out is the silent killer of joined explores: joining one-to-many without declared \texttt{relationship} inflates sums when Looker cannot apply symmetric aggregates. Declare every join's cardinality; spot-check a known total after every model change.
\end{pitfall}

\section{Wednesday: Looker Studio versus Looker, and Reverse ETL}
\index{Looker Studio vs Looker}\index{reverse ETL!identity resolution}

\subsection{Self-Serve versus Governed}
\textbf{Looker Studio} is the free, self-serve dashboard tool: connect to BigQuery (or Sheets, or CSV), drag fields, publish. It is the right answer for departmental reporting and the Thursday lab's fast path. \textbf{Looker (Core)} is the governed platform: LookML, Git workflow, access grants, API, embeds. The enterprise pattern uses both---Studio for velocity at the edges, Looker for the certified metrics at the core---with BigQuery's governed tables as the common floor.

\subsection{Reverse ETL}
Dashboards inform; reverse ETL \emph{acts}: pushing the warehouse's computed segments (high-value customers, churn risks, fraud watchlists) back into the operational systems where work happens---CRM, marketing automation, ad platforms. The reference implementation is a scheduled BigQuery query materializing the segment, plus a Cloud Function (or Dataform operation) posting deltas to the operational API on a cadence. Two disciplines make it safe: \textbf{identity resolution}---the segment keys must map to the operational system's identifiers, or you sync a segment of strangers---and \textbf{sync-frequency trade-offs}: freshness versus API rate limits and cost, which usually lands on hourly deltas, not per-event pushes.

\begin{tipbox}
Sync deltas, not full snapshots: track a \texttt{last\_synced\_at} watermark per entity and push only changed segment membership. Full syncs multiply API calls by segment size daily; delta syncs keep both cost and rate-limit pressure flat.
\end{tipbox}

\section{Thursday Hands-On Project: Streaming Dashboard and Reverse-ETL Sync}
\index{hands-on project}\index{Looker Studio dashboard}\index{Cloud Function!reverse ETL}
This lab connects Looker Studio to a BigQuery streaming-sales dataset with drill-down filters, then reverse-ETLs a high-value customer segment to a mock marketing webhook via a scheduled query and Cloud Function.

\subsection{The Live Dashboard}
\begin{enumerate}
    \item Create (or reuse) an \texttt{analytics.sales\_live} table receiving streaming inserts---the fraud capstone's alert table or a simple streaming sink works.
    \item In Looker Studio, add a BigQuery data source with a custom query selecting the trailing hour: \texttt{WHERE sale\_ts >= TIMESTAMP\_SUB(CURRENT\_TIMESTAMP(), INTERVAL 1 HOUR)}.
    \item Build tiles: revenue time series, sales-by-region map, top products table; add drill-down controls for region and product category; set data freshness to one minute.
\end{enumerate}

\subsection{The Reverse-ETL Sync}
\begin{lstlisting}[language=SQL,caption={Scheduled query materializing the high-value segment with sync watermarks.}]
CREATE OR REPLACE TABLE crm_sync.high_value_segment AS
SELECT customer_id, email, lifetime_value,
       CURRENT_TIMESTAMP() AS computed_at
FROM analytics.dim_customer
WHERE lifetime_value > 10000 AND is_current = TRUE;
\end{lstlisting}

\begin{lstlisting}[language=Python,caption={Cloud Function posting segment deltas to a mock marketing webhook.}]
import os
import requests
from google.cloud import bigquery

WEBHOOK = os.environ["MARKETING_WEBHOOK_URL"]


def sync_segment(request):
    bq = bigquery.Client()
    rows = bq.query("""
        SELECT customer_id, email FROM crm_sync.high_value_segment
        WHERE computed_at > (
          SELECT COALESCE(MAX(synced_at), '1970-01-01')
          FROM crm_sync.sync_log)
    """).result()
    payload = {"members": [dict(r) for r in rows]}
    resp = requests.post(WEBHOOK, json=payload, timeout=30)
    resp.raise_for_status()
    bq.query("INSERT INTO crm_sync.sync_log (synced_at, member_count) "
             f"VALUES (CURRENT_TIMESTAMP(), {len(payload['members'])})").result()
    return {"synced": len(payload["members"])}
\end{lstlisting}

\subsection*{Deliverables}
\begin{itemize}
    \item The dashboard URL or screenshots showing live tiles and drill-down filters.
    \item The scheduled query, Cloud Function, and sync-log table DDL.
    \item Evidence of a sync run: webhook received payload and the sync-log row.
\end{itemize}

\subsection*{Acceptance Criteria}
\begin{itemize}
    \item The dashboard reflects new streaming rows within its freshness interval.
    \item The reverse-ETL sync pushes deltas only and records each run in the sync log.
    \item The segment definition lives in exactly one place (the scheduled query), not in the Function.
\end{itemize}

\subsection*{Stretch Goals}
\begin{itemize}
    \item Rebuild one tile in Looker Core with a LookML measure and compare governance properties.
    \item Add identity resolution: map \texttt{customer\_id} to the CRM's contact IDs through a mapping table.
    \item Add alerting on sync failures and on segment-size anomalies (a segment that doubles overnight).
\end{itemize}

\section{Friday Use Case and Architecture: Governed BI for Five Hundred Analysts}
\index{governed BI}\index{metric conflicts}\index{self-service analytics}

\subsection{Scenario}
An enterprise with five hundred analysts across finance, marketing, and operations suffers metric chaos: three departments report three different ``Gross Revenue'' numbers to the same executive meeting. The mandate: one governed access layer over the central warehouse, with self-service that cannot fork definitions.

\subsection{Design}
Certify a core LookML model owned by the analytics-engineering team: conformed explores for the shared fact tables, with measures defined once, access grants per department, and content validators in CI (LookML lint plus known-total tests). Analysts build dashboards freely \emph{on the explores} but cannot redefine measures---new metric requests enter as pull requests. Looker Studio remains available for departmental scratch work against curated datasets, clearly badged as uncertified. Adoption is enforced socially as much as technically: the data-literacy program teaches ``explores are the source of truth,'' and executive reviews cite only certified dashboards.

\begin{figure}[H]
    \centering
    \resizebox{0.95\textwidth}{!}{%
    \begin{tikzpicture}[
        db/.style={cylinder, shape border rotate=90, aspect=0.25, draw=primary, thick, fill=primary!10, text width=2.9cm, minimum height=1.2cm, align=center, font=\small\bfseries},
        svc/.style={rectangle, rounded corners, draw=primary, thick, fill=primary!6, text width=3.2cm, minimum height=1.1cm, align=center, font=\small\bfseries},
        users/.style={rectangle, rounded corners, draw=codegreen!60!black, thick, fill=green!7, text width=3.0cm, minimum height=0.95cm, align=center, font=\small\bfseries},
        arrow/.style={-{Stealth[length=2.5mm]}, draw=primary, thick}
    ]
        \node[db] (bq) at (0,0) {BigQuery\\governed tables\\row policies};
        \node[svc] (lookml) at (5.0,0) {LookML model\\certified explores\\Git + CI};
        \node[users] (analysts) at (10.4,1.4) {500 analysts\\certified dashboards};
        \node[users] (studio) at (10.4,-1.4) {Looker Studio\\departmental scratch};

        \draw[arrow] (bq) -- (lookml);
        \draw[arrow] (lookml) -- (analysts);
        \draw[arrow] (bq) -- (studio);
    \end{tikzpicture}%
    }
    \caption{Governed BI: certified explores for shared truth, self-serve Studio at the edges.}
    \label{fig:ch16-governed-bi}
\end{figure}

\begin{pitfall}
A semantic layer nobody adopts is worse than none---it adds a third competing definition. Budget the socialization: training, office hours, and executive insistence on certified numbers are part of the deliverable, not an afterthought.
\end{pitfall}

\section{Interview Q\&A}
\subsection{Concepts and Trade-offs}
\begin{enumerate}
    \item \textbf{Q: Why does a LookML semantic layer prevent conflicting metric definitions?}\\
    A: Measures are defined once in versioned code and inherited by every explore, dashboard, and API call---analysts consume the definition instead of re-implementing it.

    \item \textbf{Q: What problem does reverse ETL solve that dashboards cannot?}\\
    A: Dashboards inform humans who must then act elsewhere; reverse ETL delivers the warehouse's computed decisions into the operational systems where action happens.

    \item \textbf{Q: What breaks in a sync without identity resolution?}\\
    A: Warehouse keys do not match CRM identifiers, so segments sync to wrong or nonexistent contacts---the pipeline runs green while marketing emails strangers.

    \item \textbf{Q: What trade-off governs reverse-ETL sync frequency?}\\
    A: Freshness versus API rate limits and cost; deltas on an hourly cadence are the usual optimum.

    \item \textbf{Q: When is Looker Core justified over Looker Studio?}\\
    A: When metrics must be governed---defined once, access-controlled, auditable---across many teams; Studio wins for velocity and departmental self-service.
\end{enumerate}

\section{Further Reading}
Read the Looker documentation for LookML, explores, and content governance \cite{googlecloudlookerdocs}, and Looker Studio documentation for self-serve reporting \cite{googlecloudlookerstudio}. The BigQuery documentation covers scheduled queries and the Storage Write API behind live tiles \cite{googlecloudbigquerydocs}, and Chapter~6's BI Engine reference covers dashboard acceleration \cite{googlecloudbiengine}.

\section*{Key Takeaways}
\begin{itemize}
    \item A semantic layer is code: versioned, reviewed, CI-tested---and the only durable cure for metric chaos.
    \item Declare join cardinality; fan-out inflation is silent and plausibly wrong.
    \item Live dashboards inherit warehouse performance: partitioning and BI Engine are BI decisions.
    \item Reverse ETL closes the loop from insight to action; identity resolution and delta syncs keep it safe.
    \item Studio for velocity, Looker for governance, curated tables as the common floor.
    \item Adoption is a deliverable: training and executive habits matter as much as the model.
\end{itemize}

\section{Exercises}
\begin{enumerate}
    \item \textbf{(Conceptual)} Explain symmetric aggregates to an analyst whose revenue doubled after adding a join.
    \item \textbf{(Conceptual)} Argue when \emph{not} to reverse-ETL: give two cases where a dashboard or a native integration is the better answer.
    \item \textbf{(Hands-on)} Reproduce the Thursday lab, then add a second tile showing per-city sales with a cross-filter.
    \item \textbf{(Hands-on)} Add a content-validation test to the LookML model that fails CI when \texttt{gross\_revenue} changes without a version bump in a comment convention you define.
    \item \textbf{(Design)} Write the metric-governance policy for the 500-analyst enterprise: certification tiers, PR process, SLAs for new-metric requests.
    \item \textbf{(Architecture)} Extend the reverse-ETL design with consent enforcement: only customers with marketing consent flags sync to the ad platform.
\end{enumerate}

\section{Milestone 4 Capstone: Real-Time Fraud Detection}\label{sec:ch16-capstone}
\index{Milestone 4 capstone}\index{fraud detection}\index{sliding windows!fraud}

\begin{capstonebox}
\textbf{Mission:} Simulate credit-card swipes into Pub/Sub, detect with a sliding-window Dataflow rule any card used in two different cities within ten minutes, stream alerts to BigQuery, and visualize them live in Looker---serverless end to end, with DLQ, tests, and watermark-lag alerting.

\textbf{Estimated time:} 8--10 hours.

\textbf{What you'll practice:}
\begin{itemize}
    \item Publishing simulation traffic with retry/backoff and a dead-letter topic (Chapter~13).
    \item Sliding-window stateful detection with allowed lateness (Chapter~15).
    \item Exactly-once sinks and live BI on streaming results (Chapters 14, 16).
    \item Unit-testing the fraud rule with PAssert including a late-data case.
\end{itemize}
\end{capstonebox}

\subsection*{Scenario}
Meridian Card Services wants a real-time fraud signal: a card physically presented in two distant cities within ten minutes is near-certainly compromised. Swipes stream from terminals; alerts must reach the fraud team's dashboard in under a minute; false negatives from late-arriving offline terminals are acceptable only within a documented tolerance.

\subsection*{Requirements}
\begin{itemize}
    \item Python simulator publishing plausible swipes (with seeded cross-city anomalies) to a schema-validated topic with a DLQ.
    \item Dataflow pipeline: key by card, 10-minute sliding windows with 1-minute period, two minutes allowed lateness, flag windows containing two or more distinct cities.
    \item Alerts stream to BigQuery via exactly-once write; Looker (or Studio) dashboard shows live alerts with drill-down.
    \item PAssert unit tests for the detection rule, including a late-data refinement case.
    \item Alert policies on Dataflow watermark lag and Pub/Sub backlog; documented monthly cost estimate.
\end{itemize}

\subsection*{Reference Architecture and Core Rule}
\begin{figure}[H]
    \centering
    \resizebox{0.95\textwidth}{!}{%
    \begin{tikzpicture}[
        svc/.style={rectangle, rounded corners, draw=primary, thick, fill=primary!6, text width=2.7cm, minimum height=1.0cm, align=center, font=\small\bfseries},
        warn/.style={rectangle, rounded corners, draw=red!60!black, thick, fill=red!6, text width=2.6cm, minimum height=0.9cm, align=center, font=\footnotesize\bfseries},
        arrow/.style={-{Stealth[length=2.5mm]}, draw=primary, thick},
        dasharrow/.style={-{Stealth[length=2.5mm]}, draw=codegray, thick, dashed}
    ]
        \node[svc] (prod) at (0,0) {Producers\\(card swipes)};
        \node[svc] (ps) at (3.4,0) {Pub/Sub\\(transactions)};
        \node[svc] (df) at (6.8,0) {Dataflow\\(sliding window)};
        \node[svc] (bq) at (10.2,0) {BigQuery\\(fraud alerts)};
        \node[svc] (looker) at (13.6,0) {Looker\\(live dashboard)};
        \node[warn] (dlq) at (3.4,-1.8) {Dead-Letter\\Topic};
        \node[warn] (mon) at (6.8,-1.8) {Cloud Monitoring\\(watermark lag)};

        \draw[arrow] (prod) -- (ps);
        \draw[arrow] (ps) -- (df);
        \draw[arrow] (df) -- (bq);
        \draw[arrow] (bq) -- (looker);
        \draw[dasharrow] (ps) -- (dlq);
        \draw[dasharrow] (df) -- (mon);
    \end{tikzpicture}%
    }
    \caption{Milestone 4 reference architecture: real-time fraud detection, serverless end to end.}
    \label{fig:ch16-capstone-fraud}
\end{figure}

\begin{lstlisting}[language=Python,caption={The two-cities-in-ten-minutes detection rule.}]
alerts = (
    events
    | "KeyByCard" >> beam.Map(lambda e: (e["card_id"], e))
    | "Window" >> beam.WindowInto(SlidingWindows(600, 60),
          allowed_lateness=120,
          accumulation_mode=trigger.AccumulationMode.DISCARDING)
    | "Group" >> beam.GroupByKey()
    | "Flag" >> beam.FlatMap(lambda kv: [(kv[0], w)]
        for w in [1] if len({e["city"] for e in kv[1]}) > 1))
\end{lstlisting}

\subsection*{Deliverables}
\begin{itemize}
    \item A repository that runs end-to-end from a clean environment: simulator, pipeline, tests, dashboard link.
    \item The architecture diagram including DLQ, windowing semantics, and serving layer.
    \item At least three automated tests with a failing-case demonstration.
    \item Monitoring evidence and the monthly cost estimate with two optimizations documented.
    \item A security review: least-privilege runner service account, card data pseudonymized, no hardcoded secrets.
\end{itemize}

\subsection*{Acceptance Criteria}
\begin{itemize}
    \item A seeded cross-city anomaly surfaces as a BigQuery alert row within a minute of ingestion.
    \item Late events (within tolerance) refine rather than duplicate alerts; out-of-tolerance events are measurably accounted.
    \item The pipeline survives a worker kill test without losing or duplicating alerts.
    \item The dashboard reflects the alert stream live with drill-down by city and card.
\end{itemize}

\subsection*{Stretch Goals}
\begin{itemize}
    \item Key the rule on event time from the terminal (not arrival) and quantify the difference on replayed traffic.
    \item Add a second rule---velocity (more than N swipes per minute per card)---using the State and Timers API.
    \item Reverse-ETL confirmed-fraud card IDs back to a mock authorization API to auto-decline.
    \item Load-test to ten thousand events per second and document the autoscaling curve and cost.
\end{itemize}
